% Copyright 2026 Torbjörn E. M. Nordling
% SPDX-License-Identifier: Apache-2.0
% frame-video-sync.tex — Speed from video, synchronised by the LED (16:00–18:30)
% Article §2.4 (Measurements): ellipse fit, +-23 ms quadratic window,
% \OmegaSd, \TimingSd; frame interval 1/240 s = 4.2 ms (article §4).
% Data README: sync levels; LED edges against logged times.
\begin{frame}{Measuring the speed from video}{Two clocks, one LED}
  \footnotesize
  \begin{columns}[T]
    \column{0.49\textwidth}
    \textbf{Angle and speed}
    \begin{enumerate}\setlength{\itemsep}{2pt}
      \item A tracker finds the \alert{marker} and the \alert{hub} in every frame.
      \item The marker moves on a circle that the camera sees as an \alert{ellipse}; an ellipse fit maps it back to a circle and gives the wheel angle.
      \item Each frame is timed by its own timestamp: the phone drops frames.
      \item Speed: slope of a quadratic fitted to the angle within $\pm 23$\,ms; typical $\sigma_\omega$ \OmegaSd\,rad/s.
    \end{enumerate}
    \column{0.48\textwidth}
    \textbf{Synchronising video and command log}
    \begin{enumerate}\setlength{\itemsep}{2pt}
      \item The Pi switches the LED \alert{just before each command} and logs both times on its clock.
      \item Find the LED's on and off edges in the video.
      \item Fit a straight line: video time against Pi time gives \alert{offset and drift}.
      \item Place every logged command on the video clock; timing uncertainty about \TimingSd\,ms.
    \end{enumerate}
    \vspace{2pt}
    {\scriptsize At 240 frames/s, delays below one frame interval, 4.2\,ms, are not resolved.}
  \end{columns}
\end{frame}
